%==============================================================================
% house.tex -- the house preamble for documents: articles, reports and books
%
% Loaded by the house-pdf format, after house-common.tex. Slides load
% house-beamer.tex in its place. The page itself (A4, 11pt, 2.5cm margins, one
% sided) and the fonts are format keys in _extension.yml, not LaTeX, so a
% project overrides them in YAML.
%==============================================================================

%-------------------------------------------------------------------- headings
% Headings and the title are set in the sans, Lato Semibold (the sans font's
% bold), as on a slide, so an article, a book and a deck read as one house.
% KOMA's `disposition` is sans and bold by default; \sffamily is stated so the
% house does not depend on that. Description labels are text, not headings, so
% they take the text face, bold. Guarded, because a project that overrides the
% class with a standard one (article, report, book) has no KOMA font commands.
\makeatletter
\@ifundefined{addtokomafont}{}{%
  \addtokomafont{disposition}{\sffamily}%
  \setkomafont{descriptionlabel}{\normalfont\bfseries}%
}
\makeatother

%------------------------------------------------------------- float placement
% Figures and tables stay beside the text that introduces them. greater_tables
% and pandoc both wrap a table in a bare `\begin{table}`, a float with no
% placement specifier, so LaTeX applies the class default `tbp`: top, bottom
% or its own page, but never here. The table then lands at the top of the next
% page, above the sentence that introduces it, and reads as missing. `float`
% adds the `H` specifier and `\floatplacement` makes it the default. `fig-pos:
% H` in _extension.yml does the same for figures. This is the setting the
% books used in their prefob.tex.
\usepackage{float}
\floatplacement{table}{H}
